Theories of flat-minima bias rest on mean escape-time formulas, yet the full distribution of how long an optimizer stays in a basin is rarely measured.
We measure it: we run SGD and Adam from exact minima and true saddles of small tanh networks on \fmnist and \digitsdata, across 165 conditions and 17,474 exits, and decide exit by converged quenching.
Our simulator recovers the exact Kramers rate on a double well (median ratio 1.07).
We test four numerical claims, and none holds as stated.
Exponential times are rejected in 86\% of conditions, but Weibull is the best model in only 4 of 28, its shape $k$ lies in 0.7--1.3 in only 32\%, and $k$ is unrelated to curvature ($\rho=-0.06$).
Most of the deviation from $k=1$ comes from a start-up delay (refitting after early exits gives $k=1.03$) and from mixing across minima (pooling gives $k=0.52$--$0.74$).
Doubling the batch size multiplies escape time by 5.1, not $\sqrt{2}$, while $\etaB$ still sets the time scale on \fmnist to within about 10\%.
Adam escapes about $2\times$ faster than SGD at matched in-basin loss, with no demonstrated dependence on basin width, but makes no exits at its default learning rate.
Saddle exit follows the top eigenvector 82--86\% of the time at low noise and 54\% at high noise, and the eigenvalue gap does not predict that share.
Escape is a discrete, delayed hop whose rate depends on noise and on the specific minimum.
